% Section 10 — Compute, amortization and when to search anyway (SYNTHETIC results)
\section{What a proposal costs}
\label{sec:compute}

A first proposal is cheap only after someone has paid for the searches it learned from. We therefore
charge EurekaZero for its teacher searches, every training stage and every proposal, meter all
three, never add a nested measurement twice, and convert GPU-hours to CPU-hour equivalents at a
declared rate of \syn[EZ-G18]{\EZgpuRate}, with the conclusions re-derived at half and twice that
rate in \cref{app:compute} \synbadge.

\begin{figure}[t]
\centering
% Figure — amortization: cumulative cost of one-proposal deployment against search at matched quality.
% Reads SYNTHETIC data (gap EZ-G18).
\begin{tikzpicture}
\begin{axis}[ez, width=0.62\linewidth, height=5.0cm, xmode=log, ymode=log, axis lines*=left,
    xmin=10, xmax=100000, ymin=150, ymax=3e6, clip=false,
    xlabel={problems solved (deployments)}, ylabel={cumulative CPU-hour equivalents},
    xtick={10,100,1000,10000,100000}, xticklabels={10,100,1k,10k,100k},
    ytick={1000,10000,100000,1000000}, yticklabels={1k,10k,100k,1M},
    legend pos=north west,
    legend style={font=\sffamily\fontsize{5.6pt}{6.4pt}\selectfont, draw=none, fill=white,
      fill opacity=0.85, text opacity=1, cells={anchor=west}, row sep=-1.5pt}]
  \addplot[ezgrey, line width=1pt] table[col sep=comma, x=n, y=search] {data/synthetic/cost.csv};
  \addlegendentry{search at $B_{\mathrm{eq}}$ per problem}
  \addplot[ezwarm, line width=1.2pt] table[col sep=comma, x=n, y=ez] {data/synthetic/cost.csv};
  \addlegendentry{EurekaZero: teacher search, training, then one proposal}
  \draw[ezink, line width=0.4pt, densely dotted] (axis cs:2302,150) -- (axis cs:2302,46046);
  \fill[ezink] (axis cs:2302,46046) circle (1.3pt);
  \node[font=\sffamily\fontsize{5.8pt}{6.4pt}\selectfont, text=ezink, anchor=north west, align=left]
    at (axis cs:2600,40000) {break-even\\$\approx$2,300 problems};
\end{axis}
\synfigtag
\end{tikzpicture}

\caption{Cumulative cost of solving learning-procedure problems at matched quality. Search pays
the search-equivalent budget on every problem; EurekaZero pays its teacher searches and training
once and then a negligible cost per proposal. Synthetic (gap EZ-G18).}
\label{fig:cost}
\end{figure}

\paragraph{Break-even} For learning procedures, the fixed cost is \syn[EZ-G18]{\EZfixA} CPU-hour
equivalents: \syn[EZ-G18]{\EZteachA} of teacher search and \syn[EZ-G18]{\EZshareA} for this domain's
share of training. Search that matches one proposal costs \syn[EZ-G18]{\EZsearchPerA} per problem,
and a proposal costs \syn[EZ-G18]{\EZmarginA}, so EurekaZero is cheaper after
\syn[EZ-G18]{\EZbreakA} problems (\cref{fig:cost}). For CAD, the fixed cost is
\syn[EZ-G18]{\EZfixC} against \syn[EZ-G18]{\EZsearchPerC} per specification, and break-even falls at
\syn[EZ-G18]{\EZbreakC} specifications. For next-experiment decisions the relevant cost is not
compute but experiments: a closed-loop controller that needs \syn[EZ-G04]{\EZsaveB} times fewer
evaluations to reach the target regret saves laboratory time that compute does not measure.

\paragraph{Latency} A proposal takes a median of \syn[EZ-G18]{\EZlatA}, \syn[EZ-G18]{\EZlatB} and
\syn[EZ-G18]{\EZlatC} seconds in the three games, so the decision is available while the problem is
still being discussed rather than after a search job finishes.

\paragraph{When to search anyway} The first proposal is not the best attainable answer. As for
AlphaGo Zero's raw network, which played far below the strength of the full system with search
\citep{silver2017zero}, verified search at a large enough budget overtakes one EurekaZero
proposal in both games where we measured it (\cref{fig:curves}), and the regimes of \cref{sec:heterogeneity} are where it
does so soonest. The two compose: a first proposal is a strong seed, and the search budget it saves
can be spent where the first proposal is least reliable. Search seeded by the first proposal is
reported as track S and never as a first-proposal result.

\paragraph{What we do not know about scaling} We measured one model size and four experience
fractions per game (\cref{fig:scaling}). The curves rise without saturating, but four points do not
make a scaling law, and we do not extrapolate them.
